\documentclass[11pt]{article}
\usepackage[margin=1in]{geometry}
\usepackage{amsmath,amssymb}
\usepackage{booktabs}
\usepackage[table]{xcolor}
\title{\bfseries StoryScope: what the paper actually found\\[4pt]\large A reading of Russell, Rajendhran, Pham, Iyyer \& Wieting, arXiv:2604.03136 (v6, 10 Aug 2026)}
\author{}
\date{}
\begin{document}
\maketitle
\noindent Everything below is taken from the paper itself (abstract page + full HTML). No numbers are invented. Where the paper reports a number, it is quoted; where it reports a mechanism, it is paraphrased.

\section{The question}
Most AI-text detectors look at surface style: word choice, sentence rhythm, the em-dash problem, words like ``delve'' and ``tapestry''. Those cues work, but they are fleeting. The authors point out that GPT-5.4 already cut em-dash usage sharply, and that fine-tuning an LLM to mimic human style drops detection on creative writing from 97\% to 3\% (Chakrabarty et al.\ 2026). If the surface is this easy to change, detection built on the surface will not last.

So the paper asks a different question: can a story be identified as human or AI from \emph{discourse-level narrative choices} alone---plot structure, character agency, temporal ordering, information revelation---with style stripped out entirely? Changing narrative structure is not a post-hoc edit; it is a rewrite. A detector built on narrative features should survive the edits that kill style-based detectors.

\section{The method}
\textbf{Corpus.} 10,272 human-written short stories from Books3. Each was reverse-engineered into a writing prompt (Gemini 2.5 Flash inferred the premise), and five LLMs---Gemini 3 Flash, Kimi K2.5, DeepSeek V3.2, Claude Sonnet 4.6, GPT-5.4---each wrote a story from the same prompt. Six sources total, 61,608 stories, mean 4,753 words. Source identities were anonymized in all LLM-facing prompts.

\textbf{Pipeline.} Three stages:
\begin{enumerate}
\item \emph{Structured templates.} Each story is converted by GPT-5.1 into a template organized along ten narrative dimensions from the NarraBench taxonomy: Agent, Social Network, Event, Plot, Structure, Setting, Time, Revelation, Perspective, Style. Prose becomes fields: character names and motivations, causal chains, key events, temporal order.
\item \emph{Cross-source comparison.} A held-out pool of 600 stories (100 prompts, all six versions) is presented in template form to GPT-5.1 with high reasoning effort, producing structured comparative analyses: per-source notes, cross-source divergences, recurring patterns. This compresses $\sim$2.7M tokens of raw text into $\sim$686K tokens.
\item \emph{Feature discovery.} The analyses yield 408 candidate features, each a question with typed answers (categorical, ordinal, scale, binary, multi-select). Deduplication by embedding clustering (F2LLM-4B, single linkage at cosine 0.85) gives \textbf{304 features}. A style-dependence audit excludes 47 style-linked features, leaving \textbf{257 strict narrative features} for the main experiments.
\end{enumerate}

\textbf{Classifier.} XGBoost (beat linear models and random forests). Split: 7,383 train / 1,405 val / 1,384 test prompts; tuned on validation, retrained on train+val, reported on test (8,301 stories).

\textbf{Reliability.} Production extractor (Gemini 3 Flash): Krippendorff's $\alpha=0.90$ over five runs, mean pairwise Cohen's $\kappa=0.89$. Human validation over 12 stories (240 items per annotator): human-vs-model $\kappa=0.84$, versus human-vs-human 0.74. The extractor agrees with humans at least as well as humans agree with each other.

\section{The numbers}
\textbf{Binary human-vs-AI detection (macro-F1):}
\begin{center}\small
\begin{tabular}{lccc}
\toprule
Method & Features & F1 & AUPRC \\
\midrule
Narrative (strict, no style) & 257 & \textbf{93.2} & .959 \\
Core only & 30 & 84.8 & .828 \\
Core + fingerprint & 101 & 91.1 & .934 \\
Narrative + Style & 304 & 96.0 & .982 \\
Style only & 39 & 85.8 & .867 \\
ModernBERT (raw text) & --- & 99.9 & 1.00 \\
Stylometric + XGB & 144 & 99.8 & .999 \\
TF-IDF + XGB & 5,000 & 99.7 & .999 \\
Binoculars (zero-shot) & --- & 55.9 & .404 \\
\bottomrule
\end{tabular}
\end{center}
Narrative features alone get 93.2\%---2.8 points below narrative+style, retaining over 97\% of the full model's performance. A compact 30-feature core keeps 84.8\%.

\textbf{Six-way authorship attribution (macro-F1):}
\begin{center}\small
\begin{tabular}{lc}
\toprule
Method & F1 \\
\midrule
Narrative & \textbf{68.4} \\
Core only & 46.5 \\
Core + fingerprint & 63.4 \\
Narrative + Style & 77.3 \\
Style only & 60.4 \\
\bottomrule
\end{tabular}
\end{center}
Narrative beats style-only by 8.0 points (CI 6.7--9.2). Attribution is much harder than detection; the AI models overlap heavily with each other.

\textbf{Robustness to editing.} 278 Gemini stories were run through LAMP (a span-level rewriter removing seven categories of AI artifacts, using Gemini as the rewriter). The narrative classifier still detects edited stories at \textbf{93.9\%} macro-F1 (AUPRC .988) versus 95.5\% on originals: a 1.6-point drop.

\section{What separates human from AI, in narrative terms}
\begin{itemize}
\item \textbf{AI over-explains its themes.} $\sim$20\% more explicit and moralizing on a 1--5 scale. Narrators state the theme 77\% of the time vs.\ 52\% for humans. AI dialogue serves philosophical debate 59\% vs.\ 34\%. Allusions are vague (72\% vs.\ 50\%). The meaning is said, not inferred.
\item \textbf{Human authors subvert linearity.} AI has tighter causal chains, more protagonist-driven resolutions (69\% vs.\ 46\%), fewer subplots (79\% ``no subplots'' vs.\ 57\%). AI resolutions favor internal acceptance (47\% vs.\ 27\%); humans tolerate ambiguity and use time jumps, flashbacks, nonlinear structure. A human mystery opens at the funeral and spirals backward; AI tells it from first clue to grand reveal.
\item \textbf{AI over-writes the body and senses.} Emotion through physical sensation 81\% vs.\ 38\%, smell imagery 82\% vs.\ 57\%. Humans name feelings: explicit emotion labels 29\% vs.\ 8\%. Where a human writes ``she felt afraid'', AI writes a tightening chest, cold sweat, dimming lamplight.
\item \textbf{Human authors engage the outside world.} Specific text/author references at nearly double the AI rate (47\% vs.\ 24\%), balanced explicit/implicit references (37\% vs.\ 16\%), fourth-wall breaks 67\% vs.\ 39\%, direct reader address 28\% vs.\ 7\%. AI writes as though no one is watching.
\item \textbf{AI has less diverse narrative features overall.} Humans span more locations, carry more dialogue, integrate more subplots (42\% vs.\ 21\%), give protagonists moral ambivalence more often (59\% vs.\ 38\%).
\end{itemize}

\section{The per-model fingerprints}
\begin{itemize}
\item \textbf{Claude keeps it cool.} Most distinctive profile. Restraint: flat event escalation, most uniform voice. Reverent toward literary tradition (62\% vs.\ 39--56\%). Favors epilogues, avoids dream sequences, prefers quiet endings.
\item \textbf{GPT likes to gossip.} Gossip/rumor as plot mechanism (64\% vs.\ 44--55\%), stories as reflection on events years ago, ensemble social networks. Subverts expectations more (41\% vs.\ 27--36\%), leaves reconciliations ambiguous.
\item \textbf{DeepSeek front-loads.} Puts crucial context early that other sources save.
\item \textbf{Gemini, tidiest and bleakest.} Tidiest endings, extended denouements, bleakest settings (88\% bleak and oppressive).
\item \textbf{Kimi is the generic center.} Fewest fingerprints, lowest attribution F1.
\end{itemize}

\textbf{AI convergence.} The five AI sources cluster in a shared region of narrative space, well separated from humans. Human-AI centroid distance is $1.6\times$ the AI-AI distance; even the closest human-AI pair is farther than the most distant AI-AI pair. The six most-confused attribution pairs are all AI-vs-AI; the largest (Gemini $\leftrightarrow$ DeepSeek, 222 and 207 stories) dwarfs the most common human misclassification (Human $\rightarrow$ Kimi, 46).

\textbf{Human narratives are rarer.} Rarity = mean Euclidean distance to 25 nearest neighbors. Humans: mean rarity percentile 0.71 vs.\ 0.49 (Cohen's $d=0.83$). 24.7\% of human stories in the top-10\% rarest versus 7.1\% of AI (top 1\%: 3.0\% vs.\ 0.6\%). At the prompt level, the human version is the rarest of the six 57.8\% of the time (chance: 16.7\%). Human stories are also more dispersed (22\% larger centroid radius, $1.13\times$ larger median 10-NN radius).

\section{What it means}
The durable claim is not ``we can detect AI with 93\% F1''---raw-text detectors do 99.9\%. It is that \emph{narrative structure itself carries the signal, independent of style}, and that this signal survives the edits that defeat style detectors. Changing the narrative choices means rewriting the story. So narrative features are a more durable basis for authorship analysis, and the rarity measure gives a concrete, measurable proxy for originality.

There is a flip side the paper mostly leaves implicit: the same feature space is a recipe list for making AI fiction \emph{less} detectable. Over-explain less, break the causal chain, add subplots, name real books, address the reader, let endings dangle, give the protagonist a moral mess. And since humans are rarer and more dispersed, the gap will shrink as models are pushed toward the human region of narrative space.

\section{Method caveats}
\begin{itemize}
\item Books3 short stories, mean 4,753 words. The pipeline needs long texts; transfer to flash fiction, poetry, or non-fiction is untested.
\item Human stories are one per prompt and \emph{generated} the prompt; AI stories are mirrors. The comparison is human-original vs.\ AI-mirror, not equal footing.
\item All feature extraction uses LLMs (GPT-5.1 templates, GPT-5.4 style audit, Gemini 3 Flash production). Extractor-vs-human $\kappa=0.84$ is good but not perfect; the feature space is partly an LLM's view of narrative.
\item The paper uses Books3 (copyright-problematic dataset; academic use stated). The AI disclosure notes Claude Code and Codex were used to write and polish the paper itself.
\end{itemize}

\end{document}
